\section{Exact polynomial certificates and parity identities}\label{sec:algebra}

Orient each original edge $e=uv$ arbitrarily and put
$\Delta_e=d_G(u)-d_G(v)$ and $D_e=R(u)-R(v)$. The linear form
$D_e$ does not involve $x_e$, whose contributions cancel.

\begin{proposition}\label{prop:exactpoly}
The quadratic factor
\[
 f_e=D_e^2-2\Delta_e(2-x_e)D_e
       +2\Delta_e^2(x_e-1)(x_e-2)
\]
is nonzero exactly when the actual colour condition holds on $e$,
for $x_e\in\{0,1,2\}$. Consequently $\prod_e f_e$ has a nonzero
point on that grid if and only if $\doub G$ has a legal two-colouring.
\end{proposition}
\begin{proof}
Its evaluations at $x_e=0,1,2$ are respectively
$(D_e-2\Delta_e)^2$, $D_e(D_e-2\Delta_e)$, and $D_e^2$.
These express exactly the blue-only, mixed and red-only comparisons.
\end{proof}

For arbitrary real parameters $\tau_e$ define
\[
 f_{e,\tau}=D_e(D_e-2\Delta_e)+
 \tau_e(x_e-1)(D_e+\Delta_e(x_e-2)),
\]
and the multilinear parameter polynomial
\begin{equation}\label{eq:freeC}
 C_G(\tau)=\left[\prod_e x_e^2\right]
 \prod_e(D_e^2+\tau_e x_eD_e+\tau_e\Delta_e x_e^2).
\end{equation}
The three evaluations of $f_{e,\tau}$ are
$(D_e-2\Delta_e)(D_e-\tau_e)$,
$D_e(D_e-2\Delta_e)$, and
$D_e(D_e+\tau_e-2\Delta_e)$. Nonzero therefore always implies
actual legality, although additional legal assignments may be excluded.

\begin{proposition}\label{prop:freecertificate}
If $C_G$ is not the zero polynomial, $\lir(\doub G)\le2$.
Its nonvanishing is equivalent to nonvanishing at some
$\tau\in\{0,1\}^{E(G)}$. The fixed coefficient
$c_G=C_G(2\Delta)$ is a sufficient special case.
\end{proposition}
\begin{proof}
For every polynomial of total degree at most $2m$,
\begin{equation}\label{eq:gridinterp}
 \left[\prod_e x_e^2\right]P
 =\sum_{a\in\{0,1,2\}^{E(G)}}
 \frac{P(a)}{\prod_e w(a_e)},\qquad
 w(0)=2,\ w(1)=-1,\ w(2)=2.
\end{equation}
In one variable the weighted sums of powers of degrees zero and one
vanish and that of degree two equals one. A total-degree-at-most-$2m$
monomial with any exponent above two has another below two, so only
the all-two monomial contributes. This proves the identity. Apply it
to $\prod f_{e,\tau}$: its leading coefficient is~\eqref{eq:freeC}.
A nonzero coefficient forces a nonzero grid value. The assertion
about binary parameters follows from unique interpolation of
multilinear polynomials. This is an application of~\cite{Alon99};
no general nonzero-coefficient theorem is supplied here.
\end{proof}

Write $D=Mx$. On the binary-copy variables $z_{e,1},z_{e,2}$ put
$x_e=z_{e,1}+z_{e,2}$. The copy comparison is the nonzero linear
expression $D_e-2\Delta_e+2\Delta_e z_{e,i}$. Its coefficient matrix is
\[
 \mathcal A_G=M\otimes J_2+2\operatorname{diag}(\Delta)\otimes I_2.
\]
The coefficient of the full squarefree monomial is its permanent,
and
\begin{equation}\label{eq:permanent}
 \per(\mathcal A_G)=2^m c_G.
\end{equation}
To check this equality, the quadratic part of $f_e(a+b)$ differs
from the product of the two copy linear parts only by
$2\Delta_e^2(a^2+b^2)$, which cannot contribute to a squarefree
monomial. Each substitution $x_e^2=(a+b)^2$ contributes a factor
two to the desired coefficient. Reversing an edge reverses its two
copy rows, leaving the permanent unchanged.

\subsection{A principal-minor formula modulo two}

Let $A$ and $L$ be the adjacency and Laplacian matrices of the line
graph of $G$, indexed by original edges; all determinants in the
next statement are over $\mathbb F_2$.

\begin{theorem}\label{thm:modtwo}
For every $S\subseteq E(G)$,
\[
 [\prod_{e\in S}\tau_e]C_G(\tau)
 \equiv\det L[S,S]\,\det A[\overline S,\overline S]\pmod2.
\]
Moreover, $\det A[F,F]$ is odd exactly when $(V(G),F)$ is a
forest and each of its components has odd order, including isolates.
\end{theorem}
\begin{proof}
Modulo two, $M=A$, and
$\Delta_e\equiv d_{L(G)}(e)$, so $M+\operatorname{diag}(\Delta)=L$.
After selecting the indicated parameters the product is
\[
 \prod_{e\in S}x_e(D_e+\Delta_e x_e)
 \prod_{e\notin S}D_e^2.
\]
In characteristic two the squared factors contain only squared
variables. They cannot use an $S$ variable, already present once.
Each $S$ variable still needs one odd contribution, so all the
$|S|$ linear factors must use precisely the $S$ variables, one each.
The coefficient separates into the permanents of the two principal
submatrices; modulo two these equal the stated determinants.

Let $B$ be the unsigned vertex-edge incidence matrix modulo two.
Then $A=B^TB$. If $F$ contains a cycle its incidence columns are
dependent and so is its Gram matrix. For a tree component with $p$
vertices, every $(p-1)$-row incidence minor has determinant one,
by leaf elimination. Cauchy--Binet gives its Gram determinant
$p\pmod2$. Multiply over components. Isolates contribute the
empty determinant one. This proves the forest characterisation.
\end{proof}

For an odd-order tree the constant coefficient is odd. For an
even-order tree with an even-degree vertex there is an edge $uv$
with odd $\Delta$ whose deletion gives two odd-order components:
take a leaf even-degree vertex of the smallest subtree spanning all
even-degree vertices, and an edge leaving that subtree at that
vertex. The component on the other side has only original odd degrees,
and its original degree sum is $2(|U|-1)+1$, forcing $|U|$ odd.
Taking $S=\{uv\}$ in the formula gives an odd coefficient; the
constant coefficient is even. Similarly an odd-order unicyclic graph
has an odd single-parameter coefficient if a cycle edge has
endpoints of different degree parity. These are certificates on
already known colourable graph families, not new graph-family claims.

\subsection{Offset extensions on odd forests}\label{sec:oddforest}

\begin{corollary}\label{cor:oddforest}
If each component of a forest $F$ has odd order, arbitrary real
offsets $a(v),b(v)$ admit weights $0,1,2$ such that
\[
 a(v)+\sum_{e\ni v}x_e,\qquad
 b(v)+\sum_{e\ni v}(2-x_e)
\]
both distinguish the ends of every forest edge.
\end{corollary}
\begin{proof}
The product of both colour differences has leading homogeneous part
$(-1)^m\prod_e D_e^2$. Its all-two coefficient is odd by
Theorem~\ref{thm:modtwo}, independent of offsets. Apply
\eqref{eq:gridinterp} over the reals.
\end{proof}

For physical nonnegative integer offsets with $a(v)+b(v)$ even,
this extends constructively to any odd-order connected bipartite
graph $H$. Choose the bipartition labelling $p$ or its complement
so that $\sum p\equiv\sum a\pmod2$; odd order makes the two
choices have opposite parity. Solve
$\sum_{e\ni v}x_e\equiv p(v)-a(v)\pmod2$ by a spanning tree
with non-tree weights zero. The parity sum condition makes the root
equation hold. All weights may be 0 or 1, and the final red and
blue degrees both have parity $p(v)$. Every internal edge compares
opposite parities. Neither this result nor the corollary automatically
keeps outside edges legal when boundary degrees change.

\subsection{A fixed forbidden-value application}

If each vertex of a simple graph $H$ has real offset $a(v)$ and a
fixed forbidden set $Z(v)$ with $|Z(v)|\le d_H(v)$, there is a
weighting with $a(v)+\sum x_e\notin Z(v)$ at every vertex. Enlarge
each set to size $d_H(v)$ and consider the product of all forbidden
linear factors. Its highest homogeneous part is
$\prod_v(\sum_{e\ni v}x_e)^{d_H(v)}$, with nonnegative coefficients.
Selecting each incident variable once at every vertex shows that its
all-two coefficient is at least $\prod_v d_H(v)!>0$. Interpolation
proves the assertion. More exactly that coefficient equals
\[
 \frac{\prod_v d_H(v)!}{2^{|E(H)|}}
 \#\{\text{orientations of $\doub H$ with indegree $d_H(v)$ at each $v$}\}.
\]
Indeed allocating the two occurrences of each variable to its two
endpoint factors gives the multinomial expansion; distinguishing
the two copies contributes $2!$ per edge. This also witnesses
positivity by orienting one copy each way. With blue offsets too, combine red forbidden
values with the reflected blue values; the union must still have size
at most $d_H(v)$. Final neighbour degrees that themselves depend on
these edge weights are not fixed forbidden values, so this does not
prove the conjecture.

\subsection{Limitations of the parity certificate}\label{sec:mod2limits}

If all original degrees have the same parity, $L=A$ modulo two.
A symmetric zero-diagonal matrix has every odd-order principal
determinant zero, by pairing a permutation cycle with its reversal.
Thus when $m$ is odd every coefficient of $C_G$ is even. This does
not imply zero over the integers: the triangle has $c_G=-6$.

More generally for a connected $n$-vertex graph the line adjacency
matrix has rank $n-1$ at odd $n$ and $n-2$ at even $n$. Indeed the
incidence image is the even-coordinate-sum subspace of dimension
$n-1$; the kernel of $B^T$ consists of constants, intersecting that
subspace in dimension zero or one respectively. If $m$ exceeds
twice this rank and degrees have common parity, either $S$ or its
complement is larger than the rank in every coefficient formula.
All coefficients vanish modulo two. In particular this happens for
every connected $d$-regular simple graph with $d\ge4$. Such graphs
are known to be colourable~\cite{GOW26}; the obstruction is to this
parity certificate.

Even integer nonvanishing is a substantial demand. On a regular
graph $\Delta=0$, write $Q=\prod_eD_e$. If $C_G$ is nonzero then
$Q$ has a nonzero monomial $x^\alpha$ with all $\alpha_e\le2$,
since otherwise no product with $Q$ can supply the all-two monomial.
For arbitrary real vertex offsets the same coefficient survives in
$\prod_{uv}(a(u)-a(v)+D_{uv})$. Interpolating on an arbitrary
three-element list separately assigned to every edge gives a proper
neighbour-sum weighting with those offsets. Thus this nonvanishing
would imply the corresponding $(1,3)$ list-weighting property.
Ordinary shared-list 1--2--3 weightability does not establish it.
We leave the general nonzero-polynomial question open.
